\chapter*{Answers to the exercises}\label{ch-answers}
\addcontentsline{toc}{chapter}{Answers to the exercises}

The numerical answers have been checked by running them. The checks are
programs in the directory \texttt{examples/answers} of the book's sources,
one for each chapter.

\section*{Chapter 2}\label{ans-ch2}

\begin{enumerate}
\item \code|12| and \code|1 000| are \code|int| literals (blanks are
  allowed between digits); \code|.5|,
  \code|3e2| and \code|0.0| are \code|real| literals. \code|-3| is not a
  literal but an expression: the unary operator \code|-| applied to the
  literal \code|3|. \code|1,000| is two literals, \code|1| and
  \code|000|, separated by a comma. \code|2.| is not a literal, since a
  digit must follow the point.
\item \code|17 % 5| is 3, \code|-17 % 5| is $-3$ (the unary minus binds
  first, and \code|%| truncates toward zero), \code|17 MOD 5| is 2,
  \code|-17 MOD 5| is 3 (the remainder is never negative), and
  \code|17 / 5| is the \code|real| 3.4.
\item \code|2 + 3 * 4 ** 2| is $2 + 3 \cdot 16 = 50$. \code|-3 ** 2| is
  $(-3)^2 = 9$, since unary operators bind more tightly than binary ones.
\item \code|int n = 10| binds \code|n| to the value 10; \code|n| has the
  type \code|int| and cannot be assigned to. \code|int n := 10| allocates
  a cell, stores 10 in it and binds \code|n| to the reference to the cell;
  \code|n| has the type \code|ref int|.
\item \code|ref real x = loc real := 1.5, y = loc real|. Writing
  \code|ref real| before \code|y| as well is also correct.
\item (a) is correct: \code|i| is dereferenced and the \code|int| is
  widened to \code|real|. (b) is wrong: there is no implicit conversion
  from \code|real| to \code|int|. (c) is correct. (d) is wrong:
  \code|ROUND| is defined for a \code|real| operand only, and an operand
  is not widened.
\item Both are 3. The expression \code|a + b| is evaluated first, giving
  3; this is assigned to \code|b|, and the value of \code|b| is then
  assigned to \code|a|.
\item Since a capital letter begins a type name, \code|maxCount| is read
  as the identifier \code|max| followed by the type name \code|Count|.
  Identifiers are written in lower case: \code|max_count|.
\item 2. \code|r| is a second name for the cell of \code|x|.
\item \code|((a < b) AND ((x > 0) = (y > 0))) OR (NOT done)|: the
  comparisons \code|<| and \code|>| (precedence 5) bind more tightly than
  \code|=| (4), which binds more tightly than \code|AND| (3) and
  \code|OR| (2), and the unary \code|NOT| binds most tightly of all.
\item 0.25 and 2: \code|-7 MOD 3| is 2, since \code|MOD| never yields a
  negative result. \code|NOT n = 0| means \code|(NOT n) = 0|, and
  \code|NOT| is not defined for an \code|int|.
\end{enumerate}

\section*{Chapter 3}\label{ans-ch3}

\begin{enumerate}
\item 10 and 11.
\item It prints 20 and then 1. The inner block declares a new \code|x|
  with the value $y \cdot 10 = 20$; after that block, \code|x| is the
  outermost \code|x| again.
\item (a) never, since the initial value is already greater than the
  final value; (b) five times (1, 3, 5, 7, 9); (c) five times (10, 8, 6,
  4, 2); (d) never.
\item \code|if x > 0 then "plus" elif x < 0 then "minus" else "zero" fi|.
\item \code|real|: the two branches are balanced, and the common type of
  \code|int| and \code|real| is \code|real|; the 1 is widened.
\item The operands of an operator are evaluated in an unspecified order
  (collateral elaboration), so either call may come first.
\item \code|int p := 1; while p <= 1000 do print (p); p *:= 2 od|.
\item It prints 6, 12 and 24. The number of repetitions is computed
  once, before the first iteration, so the loop runs three times, and
  the changes of \code|k| do not affect it.
\item It prints 1 forever. With a step of 0 the final value is ignored.
\item Yes. A jump is acceptable as a value of any type, here
  \code|string|. For \code|k = 5| the \code|out| part is chosen, control
  goes to \code|bad|, and \code|name| never gets a value.
\item In \code|real v = (b | x | n)| the context gives the type
  \code|real| to both branches: \code|x| is dereferenced, and \code|n| is
  dereferenced and widened. A \code|ref int| cannot become a
  \code|ref real|, so in the second declaration the \code|n| branch
  cannot be converted, and all three compilers report that
  \code|ref int| cannot be coerced to \code|ref real|.
\item The two declarations are carried out in an unspecified order, so
  \code|b| may be computed before \code|a| has a value. Written with a
  semicolon, \code|int a = 1; int b = a + 1|, the order is fixed.
\item No number is both less than 1 and greater than 9, so the
  condition is always false, and the loop ends after the first
  \code|read|. The condition must be \code|n < 1 OR n > 9|.
\end{enumerate}

\section*{Chapter 4}\label{ans-ch4}

\begin{enumerate}
\item \code|[1900:2000] real rainfall;|\\
  \code|for y from LWB rainfall to UPB rainfall do rainfall[y] := 0 od|
\item \code|a[2:3]| is the array \code|(6, 7)| with the bounds 1 to 2, so
  \code|LWB a[2:3]| is 1. \code|a[3:]| has the elements 7, 8, 9 and the
  bounds 1 to 3, so its \code|UPB| is 3. \code|a[2:3 @ 0]| has the bounds
  0 to 1, and its element 0 is 6.
\item All three are one-dimensional arrays of \code|int| (references to
  them, since \code|m| is a variable). \code|m[2, ]| is the second row,
  with the bounds 1 to 4; \code|m[, 3]| is the third column, with the
  bounds 1 to 3; \code|m[2:3, 2]| is the part of the second column in
  rows 2 and 3, renumbered to the bounds 1 to 2.
\item \code|(4, 5, 6)|: the assignment copies the elements of \code|b|
  into \code|a|, and the later change of \code|b| does not affect
  \code|a|.
\item \code|mode Date = struct (int day, month, year);|\\
  \code|Date d = (3, 11, 1975)|
\item Indexing binds more tightly than selection, so the expression means
  \code|name of (class[1][1])|, which tries to index a record. It must be
  written \code|(name of class[1])[1]|.
\item \code|y of ps := (0, 0, 0, 0, 0, 0, 0, 0, 0, 0)|. Note that
  \code|y of ps := 0| does not work: the 0 is converted into an array
  with one element, and the bounds do not match (a run-time error).
\item 2 to 11 (the second dimension is left out and keeps its bounds),
  1 to 3 (a range is renumbered), and 0 to 9 (the first dimension keeps
  its bounds).
\item \code|1 1 2 3 4 6|. The right-hand side is a value that is
  complete before anything is stored, so the overlap does no harm.
\end{enumerate}

\section*{Chapter 5}\label{ans-ch5}

\begin{enumerate}
\item (a) \code|proc (real, int) real|; (b) \code|proc bool|;
  (c) \code|proc (ref [] int) void|.
\item With a loop:
\begin{lstlisting}
proc digit_sum = (int n) int:
begin
   int sum := 0, rest := n;
   while rest > 0 do sum +:= rest MOD 10; rest := rest % 10 od;
   sum
end
\end{lstlisting}
  Recursively:
  \code{proc digit_sum = (int n) int: (n < 10 | n | n MOD 10 + digit_sum (n % 10))}.
\item The parameter \code|n| is a constant of type \code|int|, which
  cannot be assigned to. Written as \code|proc inc = (ref int n) void: n +:= 1|,
  it takes a reference to the caller's variable, as in \code|inc (count)|.
\item \code|proc apply_twice = (proc (int) int f, int x) int: f (f (x))|;
  the call \code|apply_twice ((int n) int: 2 * n, 5)| gives 20.
\item \code|prio ABOVE = 5; op ABOVE = (int a, b) bool: a > b|. Since
  \code|+| has the higher precedence 6, \code|1 + 2 ABOVE 2| means
  \code|(1 + 2) ABOVE 2|, which is \code|true|.
\item A unary operator binds more tightly than a binary one, so the
  expression means \code|(SQUARE 2) MAX 3|, which is 4.
\item An assignment is one of the expressions that are discarded as they
  are: its value, the variable \code|p|, is thrown away without calling
  the function that \code|p| holds. Every other expression that yields a
  function without parameters, such as \code|p| on its own, an element
  \code|ps[1]| or a call \code|g (1)|, is called when it is discarded.
\item 9: the precedence belongs to the symbol \code|+|, so in that block
  \code|1 + 2 * 3| is \code|(1 + 2) * 3| for integers too.
\item Yes. Two declarations of a binary operator clash only if both
  pairs of operand types are related; here the first operands
  (\code|int| and \code|real|) are not, since an \code|int| operand is
  not widened.
\item
\begin{lstlisting}
op * = ([] real a, real k) [] real:
begin
   [LWB a : UPB a] real r;
   for i from LWB a to UPB a do r[i] := a[i] * k od;
   r
end
\end{lstlisting}
\item The items of a \code|print| list are evaluated in an unspecified
  order (Section~\ref{sec-collateral}), so \code|d| may be read before or
  after \code|month| has changed it. The call must come first, in a
  declaration of its own:
  \code|string m = month (d); print ((m, " ", d))|.
\item
\begin{lstlisting}
proc sum = (int lo, hi, proc (int) real term) real:
begin
   real s := 0;
   for k from lo to hi do s +:= term (k) od;
   s
end
\end{lstlisting}
  \code|sum (1, 100, (int k) real: 1 / (k * k))| gives about 1.635, and
  \code|sum (1, 3, (int k) real: v[k] * w[k])| the scalar product of
  three-element arrays. \code|v[k] * w[k]| on its own is a \code|real|,
  not a function of \code|k|: the call is rejected, and there is no
  \code|k| at that place anyway. The function literal is called by
  \code|sum| for each \code|k|, and its body is evaluated with that
  \code|k|.
\end{enumerate}

\section*{Chapter 6}\label{ans-ch6}

\begin{enumerate}
\item $(1 + i)(1 - i) = 2$, so the real part is 2.0; \code|ABS (0 I 2)| is
  2.0.
\item There is no implicit conversion from \code|int| to \code|long int|.
  Correct is \code|long int n := long 5| or \code|long int n := LENG 5|.
\item 10, 16 and \code|true|: bit 32 is the least significant bit.
\item
\begin{lstlisting}
proc popcount = (bits b) int:
begin
   int count := 0;
   for i to bits_width do if i ELEM b then count +:= 1 fi od;
   count
end
\end{lstlisting}
\item With the environment enquiry \code|int_lengths|: \code|long long int|
  is a distinct, longer type only if \code|int_lengths| is at least 3.
\item No. The comparisons of \code|bits| are set comparisons: \code|>=|
  asks whether the first value contains all bits of the second, and
  \code|2r100| does not contain the bits of \code|2r011|.
\item \code|2 * (3 I 4)|, the complex number $6 + 8i$, since \code|I| has
  the precedence 9.
\end{enumerate}

\section*{Chapter 7}\label{ans-ch7}

\begin{enumerate}
\item 1 to 5: \code|s[2:3]| is \code|"bc"|, and the assignment gives
  \code|s| the bounds of \code|"bcxyz"|.
\item A parameter of type \code|ref [] int| would also accept a variable
  with fixed bounds, to which a longer array cannot be assigned. Only
  \code|ref flex [] int| guarantees a variable whose bounds may change.
  Moreover, in the Report a \code|ref flex [] int| is a different type,
  which cannot be passed as a \code|ref [] int| at all.
\item
\begin{lstlisting}
proc to_real = (union (int, real, bool) v) real:
   case v in (int i): i, (real r): r, (bool b): (b | 1 | 0) esac
\end{lstlisting}
\item (b) and (c). (a) is an error: the left-hand side \code|p| has the
  type \code|ref ref int|, so the right-hand side must be a
  \code|ref int|, and the compiler reports that an \code|int| cannot be
  converted to \code|ref int|.
\item One side of an identity relation is soft, the other strong.
  \code|nil| can only be the strong side, so \code|q| is the soft side and
  is not dereferenced: the variable \code|q| itself, which is never
  \code|nil|, is compared. The correct
  test is \code|ref int (q) is nil|.
\item (a) soft, (b) strong, (c) firm, (d) meek.
\item
\begin{lstlisting}
proc reverse = (ref ref Node list) void:
begin
   ref Node prev := no_node, cur := list;
   while ref Node (cur) isnt no_node do
      ref Node rest = next of cur;
      next of cur := prev;
      prev := cur;
      cur := rest
   od;
   list := prev
end
\end{lstlisting}
\item \code|words[1]| is a reference into a resizable array, a transient
  reference, and \code|+:=| is an operator whose left operand would
  receive it. \code|words[1] := words[1] + "s"| is allowed.
\item (a) is rejected: \code|"a"| is a \code|char|, which is not a member
  of the union; \code|string ("a")| works. (b) is accepted. (c) is
  rejected: \code|2| is an \code|int|, and uniting cannot follow widening.
  (d) is accepted.
\item \code|F|, \code|F|, \code|T|: \code|p is q| compares the two
  pointer variables, \code|ref int (p) is ref int (q)| compares \code|a|
  with \code|b|, and \code|ref int (p) is a| compares \code|a| with
  \code|a|. After \code|q := a| the results are \code|F|, \code|T|,
  \code|T|: the variables \code|p| and \code|q| are still different
  cells.
\item With \code|push|, the array is replaced when it holds 4, 8, 16,
  \ldots, 512 elements, which copies $4 + 8 + \cdots + 512 = 1020$
  elements, about one per element. With \code|append|, the $k$-th call
  copies the $k - 1$ elements before it, $0 + 1 + \cdots + 999 = 499\,500$
  in all.
\item The function \code|(int x) int: x + n| uses \code|n|, a parameter
  of \code|adder|, so its lifetime ends when \code|adder| returns, and it
  may not be returned. Algol 68 Genie stops with ``value is exported out
  of its scope''; ga68 and a68j run the program, and the
  function returned by \code|adder (3)| even yields 7 for the argument 4.
\item The first loop body declares nothing, since \code|print (i)| is not
  a declaration: the program remains correct. The second declares
  \code|j|, so each cell lives only during one execution of the body, and
  the program is wrong; Algol 68 Genie stops it.
\item \code|a is b| is \code|false|: the two lists consist of different
  cells, although their contents are equal.
\begin{lstlisting}
prio EQ = 4;
op EQ = (ref Node a, b) bool:
   if a is nil then b is nil
   elif b is nil then false
   else value of a = value of b andth (next of a EQ next of b)
   fi
\end{lstlisting}
  \code|a EQ b| is \code|true|.
\end{enumerate}

\section*{Chapter 8}\label{ans-ch8}

\begin{enumerate}
\item Both leave the order of \code|a| and \code|b| unspecified, and in
  both they may be interleaved. With \code|par|, each is a separate
  process, so one of them can wait on a semaphore while the other goes on.
\item The part waits forever; if no other part can continue either, the
  program hangs (a deadlock). Formally, the Report leaves the rest of the
  program undefined once no process can continue.
\item With one producer and one consumer, \code|mutex| is not strictly
  necessary: \code|tail| is changed only by the producer and \code|head|
  only by the consumer, and \code|free| and \code|full| ensure that they
  never use the same element of the buffer at the same time. With several
  producers or several consumers, two of them could change the same
  position at once, and \code|mutex| becomes necessary.
\item Let philosopher 0 pick up the right fork first and then the left
  one, while all others pick up the left fork first. Then the five can
  never each hold exactly one fork while waiting for the next.
\item When the buffer is full, the producer takes \code|mutex| and then
  waits in \code|DOWN free| while holding it. The consumer waits in
  \code|DOWN mutex| and can never make a place free: a deadlock.
\item \code|sema| is a record whose only field is a reference to the
  counter, so \code|b| is a copy of that reference and refers to the
  same counter as \code|a|.
\end{enumerate}

\section*{Chapter 9}\label{ans-ch9}

\begin{enumerate}
\item \code|open| connects the file with an existing book;
  \code|establish| creates a new one.
\item That the routine has dealt with the event, so that the transput
  continues at the same place. If it returns \code|false|, the default
  action follows; at the end of the file, the Report leaves the program
  undefined, and the implementations stop it with an error.
\item
\begin{lstlisting}
proc int_to_string = (int n) string:
begin
   string s := "";
   File f;
   associate (f, s);
   put (f, n);
   s
end
\end{lstlisting}
  The result has the standard layout of \code|print|, with the sign and
  leading blanks.
\item \code|$4d$| prints \texttt{0007} for 7; for $-7$ it is an error,
  since the picture has no sign frame. \code|$4zd$| prints four blanks
  and 7 for 7, and is also an error for $-7$. \code|$-3zd$| prints four
  blanks and 7 for 7, and a minus sign, three blanks and 7 for $-7$: the
  sign frame \code|-| prints a blank for a positive value.
\item Reading a string stops before the line break and does not skip it,
  so the line break must be skipped with \code|newline|. If both are read
  in one call, \code|get (input, (line, newline))|, and the last line has
  no line break, \code|newline| meets the end of the input and the event
  routine jumps away before the line is printed. With two calls, the line
  is printed first.
\item (a) One value: the frames form a single picture of four digits,
  and 3 is printed as \code|00/03|; the format then starts again and
  prints 4 as \code|00/04|. (b) Two values, printed as
  \code|03/04|.
\item \code|" -5"|, \code|"  5"|, \code|"2.50"| and \code|"-.25"|.
  A positive width always writes a sign, a negative one only a minus
  sign; with width 0, \code|fixed| writes no zero before the point.
\item Both \code|a| and \code|b| fail on the same \code|x|, since the
  routine leaves it in place; in the implementations, the routine is
  called twice, and \code|a| and \code|b| keep their old values. The
  \code|5| is never read.
\item
\begin{lstlisting}
$g(0), " ", c("January", "February", "March", "April", "May", "June",
   "July", "August", "September", "October", "November", "December"),
 " ", g(0)$
\end{lstlisting}
  The pictures \code|" "| have no pattern and use no value.
\end{enumerate}
